\documentclass[xcolor=table]{beamer}
\usepackage{beamerthemesplit}
\usepackage{wrapfig}
\usetheme{SPbGU}
\usepackage{pdfpages}
\usepackage{amsmath}
\usepackage{cmap}
\usepackage[T2A]{fontenc}
\usepackage[utf8]{inputenc}
\usepackage[english]{babel}
\usepackage{indentfirst}
\usepackage{amsmath}
\usepackage{tikz}
\usepackage{multirow}
\usepackage[noend]{algpseudocode}
\usepackage{algorithm}
\usepackage{algorithmicx}
\usepackage{fancyvrb}
%usepackage{fancyvrb}
%\usepackage{minted}
%\usepackage{verbments}

\newenvironment{psmallmatrix}
  {\left(\begin{smallmatrix}}
  {\end{smallmatrix}\right)}


\beamertemplatenavigationsymbolsempty

\title[SCFPQ]{Синтаксический анализ графов для случая стохастических грамматик и вероятностных графов с помощью систем матричных уравнений}
\institute[СПбГУ]{
JetBrains Research, Programming Languages and Tools Lab  \\
Санкт-Петербургский Государственный Университет
}

\author[Юлия Сусанина]{Юлия Сусанина}

\date{19.12.2020}

\begin{document}
{
\begin{frame}[fragile]
  \begin{tabular}{p{2.0cm} p{7.5cm} p{1cm}}
   \begin{center}
      \includegraphics[height=1.5cm]{pictures/jetbrainsResearch.pdf}
    \end{center}
    &
    \begin{center}
      \includegraphics[height=1.5cm]{pictures/YC_logo.pdf}
    \end{center}
    &
    \begin{center}
      \includegraphics[height=1.5cm]{pictures/SPbGU_Logo.png}
    \end{center}
  \end{tabular}
  \titlepage
\end{frame}
}

\begin{frame}[fragile]
  \transwipe[direction=90]
  \frametitle{Context-Free Path Querying (CFPQ)}
  
\begin{itemize}
  \item КС-грамматика $G=(N, \Sigma, R)$ \\
  $\mathcal{L}(G_S) = \{\omega \mid S \Rightarrow_{G}^{*} \omega\}$, $S \in N$
  \item Орграф $D = (V,E,\sigma)$, $\sigma \subseteq \Sigma$, $E\subseteq V\times \sigma \times V$ \\
  $m\lambda n$ --- путь из $m$ в $n$ в графе $D$, $\lambda$ --- слово данного пути
  \item $R_A = \{(m, n) \mid m\lambda n \text{ --- путь в } D, \lambda \in \mathcal{L}(G_A)\}$
\end{itemize}
\end{frame}

\begin{frame}[fragile]
  \transwipe[direction=90]
  \frametitle{CFPQ на языке линейной алгебры}
\begin{itemize}
  \item Через перемножение булевых матриц
            \begin{minipage}{.9\linewidth}
            \begin{algorithm}[H]
            \begin{algorithmic}
                \ForAll{$(i,x,j) \in E$}
                    \State{$T_{i,j} \gets T_{i,j} \cup \{A~|~(A \rightarrow x) \in P \}$}
                \EndFor
                \While{матрица $T$ меняется}
                    \State{$T \gets T \cup (T \times T)$}
                \EndWhile
            \end{algorithmic}
            \end{algorithm}
            \end{minipage}
  \item Как системы матричных уравнений над $\mathbb{R}$
            \begin{minipage}{.9\linewidth}
            \begin{algorithm}[H]
            \begin{algorithmic}
            \ForAll {$N_i \to \beta^0_0 \dots \beta^0_k \mid \ldots \mid \beta^l_0 \dots \beta^l_m$}
                \State{решить $\mathcal{T}_{N_i} = \epsilon_{N_i}(T_{\beta^0_0}\cdot \ldots \cdot T_{\beta^0_k} + \ldots + T_{\beta^l_0}\cdot \ldots \cdot T_{\beta^l_m}),$}
                \State{$ \text{где } \epsilon_{N_i} \text{выбрана так, чтобы} \mathcal{T}_{N_i}^{k} \leq \textbf{1}$}
            \EndFor
            \end{algorithmic}
            \end{algorithm}
            \end{minipage}
\end{itemize}

\end{frame}



\begin{frame}[fragile]
  \transwipe[direction=90]
  \frametitle{Мотивация}
\begin{itemize}
  \item Данные из реального мира почти всегда могут быть получены только с некоторой точностью. Детерминированные модели нередко слишком грубо описывают действительность
  \item Данные могут содержать ошибки и неточности
\end{itemize}

$$\Downarrow$$

\begin{center}
    Необходимо находить новые вероятностные модели для более точной обработки существующих массивов данных (например, из биоинформатики)
\end{center}
\end{frame}

\begin{frame}[fragile]
  \transwipe[direction=90]
  \frametitle{Stochastic Context-Free Path Querying (SCFPQ)}
  \begin{itemize}
  \item Пути с наибольшей вероятностью для всех пар вершин
      $$T_A[i,j] == \max_{\substack{A \to BC \\ 0 \leq k < n}} \Theta(A \to BC) T_B[i,k] T_C[k,j]$$
  \begin{itemize}
      \item Сложность: $(|N^3||V^5|)$
  \end{itemize}
  \item Вероятности между всеми парами вершин
      \begin{align*}
        T_A[i,j] &= T_A[i,j] + \Theta(A \to BC) \ (T_BT_C)[i, j] \\
        \text{or} \\
        T_{N_i} &= \sum_{N_i \to \alpha_j} \Theta(N_i \to \alpha_j) \ T_{\beta^j_0}\cdot \ldots \cdot T_{\beta^j_k}
    \end{align*}
  \begin{itemize}
      \item В общем случае точное вычисление невозможно
  \end{itemize}
    
\end{itemize}
\end{frame}


\begin{frame}[fragile]
  \transwipe[direction=90]
  \frametitle{Пример}
  
\begin{columns}
\begin{column}{0.6\textwidth}
  
$$S \xrightarrow[]{0.3} aSb$$
$$S \xrightarrow[]{0.6} cSd$$

\only<1>{
\def\a{\color{red}0.03}
\def\b{\color{red}0.00081}
\def\c{\color{red}0.0027}
\def\d{\color{red}0.1}
\def\e{\color{red}0.000243}
\def\f{\color{red}0.009}
\def\g{\color{violet}0.1}
\def\k{\color{violet}0.036}
\def\l{\color{violet}0.01296}
\def\n{\color{violet}0.0216}
\def\m{\color{violet}0.07776}
\def\o{\color{violet}0.06}

$$$$

    \begin{psmallmatrix}
    0&0&\a& &\b&0& \\   
    0&0&\c&&\d&0& \\
    0&0&\e&&\f&0&  \\ 
    0&0&0&0&0&0&0 \\
    0&\g&0&0&0&\k&\l \\
    0&\n&0&0&0&\m&\o    \\
    0&0&0&0&0&0&0      \\
    \end{psmallmatrix}
}

\only<2>{
\def\a{\color{red}30000}
\def\b{\color{red}810}
\def\c{\color{red}2700}
\def\d{\color{red}100000}
\def\e{\color{red}243}
\def\f{\color{red}9000}
\def\g{\color{violet}100000}
\def\k{\color{violet}36000}
\def\l{\color{violet}12960}
\def\n{\color{violet}21600}
\def\m{\color{violet}7776}
\def\o{\color{violet}60000}

$$$$

    \begin{psmallmatrix}
    0&0&\a& &\b&0& \\   
    0&0&\c&&\d&0& \\
    0&0&\e&&\f&0&  \\ 
    0&0&0&0&0&0&0 \\
    0&\g&0&0&0&\k&\l \\
    0&\n&0&0&0&\m&\o    \\
    0&0&0&0&0&0&0      \\
    \end{psmallmatrix}
}
    
\end{column}
\begin{column}{0.4\textwidth}  %%<--- here

    \begin{tikzpicture}[->,node distance={22mm}, main/.style = {draw, circle}] 
    
    \node[main] (5) {$x_6$};
    \node[main] (1) [below left of=5] {$x_2$};
    \node[main] (0) [below left of=1] {$x_1$};
    \node[main] (2) [below right of=1] {$x_3$}; 
    \node[main] (3) [below left of=2] {$x_4$};
    \node[main] (6) [above left of=1] {$x_7$};
    \node[main] (4) [above left of=5] {$x_5$};
    
      \path
        (0) edge[red] [above] node [pos=0.4] {a} (1)
        (1) edge[red] [above] node [pos=0.6]  {a} (2)
        (2) edge[red] [above] node {a} (0)
        (2) edge[red] [bend right] node [above, pos=0.6] {b} (3)
        (3) edge[red] [bend right] node [below, pos=0.6] {b} (2)
        (4) edge[violet] [bend left] node [above, pos=0.6] {c} (5)
        (5) edge[violet] [bend left ] node [below, pos=0.6] {c} (4)
        (1) edge[violet] [below] node [pos=0.6] {d} (6)
        (6) edge[violet] [below] node {d} (5)
        (5) edge[violet] [below] node [pos=0.4]  {d} (1);
    \end{tikzpicture}

\end{column}
\end{columns}

\end{frame}


\begin{frame}[fragile]
  \transwipe[direction=90]
  \frametitle{Взвешенные правила и терминалы}
\begin{itemize}
  \item Веса у правил грамматики --- приоритет путям с определенными паттернами
  \item Веса у вершин графа (или терминалов) --- приоритет путям, проходящим через определенные ребра
  \item Иногда эти два случая совпадают
  \begin{itemize}
      \item Веса у правил грамматики и у терминалов для однозначных грамматик
  \end{itemize}
\end{itemize}
\end{frame}



\begin{frame}[fragile]
  \transwipe[direction=90]
  \frametitle{Результаты}
  
\begin{itemize}
    \item Сформулированы две проблемы синтаксического анализа графов для случая стохастических грамматик и вероятностных графов
    \item Предложены методы решения этих проблем с помощью методов, основанных на алгоритмах линейной алгебры и вычислительной математики
    \item В процессе
    \begin{itemize}
    \item Создание эталонного набора данных
    \item Эффективная параллельная реализация предложенных алгоритмов с помощью разных итеративных методов и их сравнение
    \item Короткая статья на EDBT 2021
\end{itemize}
\end{itemize}
\end{frame}


\begin{frame}[fragile]
  \transwipe[direction=90]
  \frametitle{Публикации прошлых работ}
  
\begin{itemize}

  \item \textbf{ACM SIGMOD 2020 Student Research Competition:} \\ 
  Yuliya Susanina. Context-Free Path Querying via Matrix Equations.\\
  \item \textbf{LNBI 2020:} \\
  Yuliya Susanina, Anna Yaveyn and Semyon Grigorev. Modification of Valiant’s Parsing Algorithm for String-Searching Problem. \\
  \item \textbf{Журнал <<Труды ИСП РАН>>:} \\
  Сусанина Ю.А., Григорьев С.В. Модификация алгоритма Валианта для задачи поиска подстрок. \\
  
\end{itemize}

\end{frame}

\end{document}